% A math field that is neither a character nor a braced list.
%
% Measured against pdftex: the reference puts a { in front of it, so the
% field becomes a list of its own, and says so --
%
%   ! Missing { inserted.
%   A left brace was mandatory here, so I've put one in.
%   ...
%
% -- and the } it wants turns up later, when the formula or the group
% above it ends:
%
%   {math mode: \mathord}
%   ! Missing { inserted.
%   {\radical}
%   {math shift character $}
%   ! Missing } inserted.
%
% So `$\mathord\radical"161 x$' draws BOTH, and the box that comes out is
% the same one this engine already builds in silence.
%
% Which commands can stand as a field without the brace: a character, and
% \char, \mathchar and \delimiter (and a { of course). Everything else
% draws the pair -- \radical, \hbox, \vcenter, \mathop were all measured.
% \relax is skipped before the field is looked for and draws nothing.
%
% WHERE THIS ENGINE STANDS: it builds the right box and says nothing,
% except for a field-taking command met while a field is already expected
% -- `$\underline\radical"161 x$' -- which stops the run.
%
% NOT YET. hstex has no scan-a-math-field step: \mathord sets a forced
% class and \underline sets a slot for whatever comes next to fill, so
% there is no single place the brace would be inserted. What makes this
% worth doing is that the closing half is already right: hstex reports
% "Missing } inserted" for a group left open in a formula, so opening a
% real math group at the right moment would give both messages. What it
% needs is a list of the commands that may stand as a field, and getting
% that list wrong moves boxes rather than diagnostics.
%
% trip draws this on lines 272, 396, 402 and 409 among others.
\tracingonline=1\showboxbreadth=9999 \showboxdepth=9999
\font\tt=cmr10 \tt
\setbox0=\hbox{$\mathord\radical"161 x$}\showbox0
\setbox1=\hbox{$\mathord\hbox{Q}$}\showbox1
\setbox2=\hbox{$\mathord\mathchar"161 $}\showbox2
\setbox3=\hbox{$\mathord X$}\showbox3
\setbox4=\hbox{$\mathord\relax x$}\showbox4
\end
